% ─── Ключевые выводы ─────────────────────────────────────────────────────
\begin{frame}[shrink=15]{Ключевые выводы}

  \begin{enumerate}\small
    \setlength\itemsep{0.15cm}

    \item \textbf{LSTM --- лидер при честном сравнении.} В условиях однократного обучения нейросеть LSTM (MAPE = 0,33\%) значительно превосходит Random Forest (0,88\%) и статическую ARIMA (1,89\%).

    \item \textbf{Rolling forecast --- мощный, но методологически особый.} ARIMA с rolling forecast (MAPE = 0,15\%) показывает наименьшие ошибки, однако это преимущество обусловлено постоянным переобучением, а не качеством самой модели.

    \item \textbf{Нейросети лучше улавливают нелинейность.} LSTM способна обучиться сложным паттернам валютного ряда, тогда как статическая ARIMA теряет точность уже на среднесрочном горизонте.

    \item \textbf{Random Forest уступает обоим подходам.} Несмотря на устойчивость к шуму, ансамблевый метод не смог эффективно моделировать автокорреляционную структуру ряда USD/UZS.

    \item \textbf{Практический вывод.} Для оперативного прогнозирования с обновлением данных оптимальна ARIMA (rolling). Для автономных прогнозов без переобучения --- LSTM.
  \end{enumerate}
\end{frame}
